\section{Certificates as deficit bounds}\label{sec:deficit}

The obstruction certificates of the calculus prove that an information state is nonviable. Read through the safety value, each certificate is a computable lower bound on the exact deficit, and its attaining witness is the dual object of the value---a lower-bound side of the exact number, not a surrogate for it.

\begin{proposition}[deficit lower bounds]\label{prop:deficit}
Write $\Delta_k(b)=1-V_k(b)$ and let $p(x,a)$ be the one-step exit probability of Definition~\ref{def:model}.
\begin{enumerate}[leftmargin=1.6em,itemsep=1pt]
\item[(i)] \emph{Chance-constrained common action.} If $\sum_{x}b(x)\,p(x,a)\ge\delta$ for every $a\in A$, then $\Delta_k(b)\ge\delta$ for every $k\ge1$. This holds for arbitrary (stochastic) kernels and observation laws.
\item[(ii)] \emph{Minimal mass.} Under deterministic kernels, if no admissible policy of the class keeps the whole support alive for $k$ stages, then $\Delta^{\Pi}_k(b)\ge\min_{x\in\supp b}b(x)$; the exact deficit is $\min_S b(S^{c})$ over the maximal survivable sets (Theorem~\ref{thm:antichain}(iv)).
\item[(iii)] \emph{Monotonicity.} $\Delta_{k+1}(b)\ge\Delta_k(b)$; the deficit is non-decreasing in the horizon and freezes after finitely many strict increases (Proposition~\ref{prop:freeze}).
\end{enumerate}
\end{proposition}

\begin{proof}
(i) One-step survival under $a$ has probability $\sum_x b(x)(1-p(x,a))\le1-\delta$, uniformly in $a$, so $V_1(b)\le1-\delta$, and Theorem~\ref{thm:support}(iii) (whose proof uses only absorption) extends the bound to every $k\ge1$. (ii) is Theorem~\ref{thm:antichain}(iv): every maximal survivable set misses a state of mass at least $\min b$. (iii) is Theorem~\ref{thm:support}(iii) and Proposition~\ref{prop:freeze}.
\end{proof}

Two instances shared with the calculus show the bounds attained.

\begin{example}[two-floor instance]\label{ex:twofloor}
States $\{0,1\}$, actions $u\in\{\tfr25,\tfr35\}$, no informative observation; state $0$ survives exactly the lower action and state $1$ exactly the higher, each indefinitely. The one-step witness set is $\Gamma_1=\{(1,0),(0,1)\}$ and $V_k(b)=\max(b_0,b_1)$ for every $k\ge1$ (Figure~\ref{fig:deficits}a). At the symmetric prior the chance-constrained certificate fires with $\delta=\tfr12$, the exact deficit is $\tfr12$, and it equals the minimal mass: both bounds of Proposition~\ref{prop:deficit} are attained simultaneously.
\end{example}

\begin{example}[asymmetric audit on the delayed instance]\label{ex:asym}
On the delayed hidden-regime instance of Section~\ref{sec:delayed} at $(z_0,\Tobs)=(\tfr32,2)$ with prior mass $w=\tfr3{10}$ on the regime $\theta=+1$, the four blind two-step sequences keep alive the masses $\tfr3{10},\tfr3{10},\tfr7{10},\tfr7{10}$: every sequence loses a branch at its first step, because at $z_0=\tfr32$ the first move sends one branch to $\tfr12<1$ (Table~S2.1). Hence $V_2(b_0)=\tfr7{10}$ and the deficit $\tfr3{10}$ equals the minimal mass---the bound \eqref{eq:minmass} is attained away from the symmetric prior, by the maximal survivable set that misses exactly the lighter branch.
\end{example}

\begin{figure}[t]
\centering
\includegraphics[width=\linewidth]{fig1_deficits.pdf}
\caption{Deficits attained and deficits below the minimal mass. (a) The two-floor instance: the value is $\max(b_0,b_1)$ for every horizon, the two witnesses $(1,0)$ and $(0,1)$ are the indicators of the two maximal survivable singletons, and the shaded deficit equals the minimal mass at every prior. (b) The two-state noisy sensor of Proposition~\ref{prop:pathdeficit}: the optimized one-step deficit is $\min\{\varepsilon,b_1,b_2\}$, which falls strictly below the minimal mass (dashed) on the whole interior plateau; at $b=(\tfr12,\tfr12)$ and $\varepsilon=\tfr1{20}$ the deficit is $\tfr1{20}$ against a minimal mass of $\tfr12$.}
\label{fig:deficits}
\end{figure}
